\scanpage{101}
\mspage{100}
\mspage{22.}
proves one half of (6.5) in all cases.

The problem of determining the actual numbers $q_k$ in (6.5) is thus
seen to be a special case of the more general problem of determining
the structure of the group $J(X)$. An attack on this problem by the
methods of [2] and [3] will be given in a future publication.

We conclude with a few remarks on the relation between the methods of
this paper and those of James. The main tool of James is the ``intrinsic
join''. Actually he defines several ``intrinsic joins'', two of them
being maps
\[
\begin{aligned}
Q_{n,k}*P_{m,k}&\longrightarrow P_{m+n,k}\\
Q_{n,k}*Q_{m,k}&\longrightarrow Q_{m+n,k},
\end{aligned}
\]
where $X*Y$ denotes the join of $X$ and $Y$. Since
$X*Y=S(X\mathbin{\times\mkern-10mu\times}Y)$\inserted{${}^+$}
these two maps can, in view of (2.4) (4.3) and (5.4), be written
\[
\begin{aligned}
P_k^{(n-k)\xi\oplus\eta}\mathbin{\times\mkern-10mu\times}
P_k^{(m-k)\xi}&\longrightarrow P_k^{(m+n-k)\xi}\\
P_k^{(n-k)\xi\oplus\eta}\mathbin{\times\mkern-10mu\times}
P_k^{(m-k)\xi\unsure{\oplus\eta}}&\longrightarrow
P_k^{(m+n-k)\xi\oplus\unsure{\zeta|\eta}}.
\end{aligned}
\]
In view of (4.5) both these are maps of the type
\[
f:X^\alpha\mathbin{\times\mkern-10mu\times}X^\beta
\longrightarrow X^{\alpha+\beta+\tau}
\]
where $X$ is a compact differentiable manifold with tangent bundle
$\tau$. By (3.3), and the fact that the smash product commutes with
duality, the dual of $f$ is a map of the form:
\[
Df:X^{\gamma+\delta}\longrightarrow
X^\gamma\mathbin{\times\mkern-10mu\times}X^\delta
\]
where $\gamma=-\alpha-\tau$, $\delta=-\beta-\tau$. Now the diagonal map
$X\longrightarrow X\times X$ induces a map
\[
\Delta:X^{\gamma+\delta}\longrightarrow(X\times X)^{\gamma\times\delta}
=X^\gamma\mathbin{\times\mkern-10mu\times}X^\delta
\quad\text{by (2.3).}
\]
It would seem reasonable to suppose that the maps $Df$ and $\Delta$ are
of the same $S$-homotopy type.

In addition to the maps first mentioned James also defines an intrinsic
join for the Stiefel manifolds $O_{n,k}$. This has no counter-part here,
since we are not directly concerned with Stiefel manifolds.

\margin{$+$ up to homotopy type}
Pembroke College,\\
Cambridge.

